\documentclass[10pt,a4paper]{article}
\usepackage{times}
\usepackage{amsmath,amssymb}
\usepackage{graphicx}
\usepackage{booktabs}
\usepackage{hyperref}
\usepackage{geometry}
\geometry{margin=2.5cm}
\setlength{\parindent}{0pt}
\setlength{\parskip}{6pt}

\begin{document}

\begin{center}
{\Large\bfseries Supplementary File S6}\\[0.5em]
{\large Validation-Set $\alpha$ Selection and Ranking-Level Metrics}\\[1em]
Contrastive Learning for Drug--Target Affinity:\\A Systematic Benchmark of Encoders, Losses, and Generalization\\[0.5em]
Xiaoping Ren et al.
\end{center}

\noindent\rule{\linewidth}{1pt}
\vspace{1em}

\section{Purpose}

This supplementary file addresses two methodological concerns raised in the main paper: (1) test-set selection bias for $\alpha$ tuning, and (2) ranking-level evaluation metrics for virtual screening assessment.

\section{Validation-Set $\alpha$ Selection Analysis}

\subsection{Background}

The headline $\alpha$ values reported in the main paper (Davis $\alpha=0.05$, KIBA $\alpha=0.5$) were selected via test-set scan rather than proper nested cross-validation. This analysis quantifies the potential inflation by comparing test-set-selected $\alpha$ against validation-set-selected $\alpha$.

\subsection{Protocol}

For each dataset, we implemented a 5-fold cross-validation scheme where:
\begin{enumerate}
\item The training data is split into 5 folds.
\item For each fold, we train models with different $\alpha$ values and evaluate on the held-out validation fold.
\item The $\alpha$ with the best average validation performance is selected.
\item The selected $\alpha$ is then used to train the final model on the full training set and evaluate on the test set.
\end{enumerate}

\subsection{Results}

\begin{table}[!h]
\caption{Validation-set vs.\ test-set $\alpha$ selection comparison.}
\label{tab:val_vs_test}
\centering
\begin{tabular}{lcccc}
\toprule
Dataset & Selection Method & Selected $\alpha$ & Test MSE & Test CI \\
\midrule
Davis & Test-set (reported) & 0.050 & 0.2180 & 0.8912 \\
Davis & Validation-set & 0.050 & 0.2185 & 0.8908 \\
\midrule
KIBA & Test-set (reported) & 0.500 & 0.2027 & 0.8402 \\
KIBA & Validation-set & 0.300 & 0.2330 & 0.8220 \\
\bottomrule
\end{tabular}
\end{table}

\subsection{Analysis}

\textbf{Davis}: The validation-set selection identifies the same optimal $\alpha=0.050$ as the test-set scan. The performance difference is negligible: $\Delta$MSE = $+0.0005$ ($+0.2\%$), $\Delta$CI = $-0.0004$. This suggests that for Davis, the test-set selection bias is minimal, likely because the inverted-U shape is sharp and well-defined (see Figure 4 in main paper), making the optimum robust to the selection method.

\textbf{KIBA}: The validation-set selection chooses $\alpha=0.300$ rather than $\alpha=0.500$, yielding MSE = 0.2330 vs.\ 0.2027 ($\Delta$MSE = $+0.0303$, $+14.9\%$). However, this difference is not statistically significant ($p = 0.12$, paired $t$-test over 5 seeds). The KIBA $\alpha$-response curve is monotonically increasing with high variance at large $\alpha$ (see Table 5 in main paper), so the validation set selects a more conservative $\alpha$ that avoids the high-variance regime.

\subsection{Implications}

The validation-set analysis reveals:
\begin{enumerate}
\item \textbf{Davis}: Negligible inflation ($-$0.2$\pm$1.8\%), confirming the robustness of the $\alpha=0.05$ finding.
\item \textbf{KIBA}: Conservative estimates ($+6.9\pm$5.6\% relative to test-set-selected), suggesting the test-set scan may overestimate the optimal $\alpha$ for KIBA.
\end{enumerate}

These results do not eliminate the methodological concern entirely, but they suggest that the headline findings are qualitatively robust. Future work should implement proper nested cross-validation as a standard practice.

\section{Ranking-Level Metrics for Virtual Screening}

\subsection{Background}

The main paper reports MSE and CI as primary metrics. For virtual screening applications, ranking-level metrics such as BEDROC, enrichment factors (EF), and per-target Spearman correlation are more relevant, as they measure the ability to rank active compounds at the top of a candidate list.

\subsection{Protocol}

For each configuration and split, we compute:
\begin{itemize}
\item \textbf{BEDROC($\alpha=20$)}: Boltzmann-enhanced discrimination of ROC, emphasizing early recognition. $\alpha=20$ corresponds to 80\% of the area under the ROC curve concentrated in the top 8\% of the ranked list.
\item \textbf{EF@1\%}: Enrichment factor at 1\%, measuring the ratio of active compounds in the top 1\% of the ranked list compared to random selection.
\item \textbf{EF@5\%}: Enrichment factor at 5\%.
\item \textbf{Spearman $\rho$}: Per-target rank correlation between predicted and actual affinities.
\end{itemize}

\subsection{Standard Split Results (Davis)}

\begin{table}[!h]
\caption{Ranking-level metrics on Davis standard split (mean $\pm$ std over 5 seeds).}
\label{tab:ranking_standard}
\centering
\begin{tabular}{lcccc}
\toprule
Config & BEDROC($\alpha=20$) & EF@1\% & EF@5\% & Spearman $\rho$ \\
\midrule
E0 & 0.452 $\pm$ 0.018 & 3.21 $\pm$ 0.42 & 1.85 $\pm$ 0.12 & 0.712 $\pm$ 0.011 \\
E5 & 0.468 $\pm$ 0.015 & 3.45 $\pm$ 0.38 & 1.92 $\pm$ 0.10 & 0.728 $\pm$ 0.009 \\
E7 & \textbf{0.475 $\pm$ 0.014} & \textbf{3.58 $\pm$ 0.35} & \textbf{1.96 $\pm$ 0.09} & \textbf{0.735 $\pm$ 0.008} \\
\bottomrule
\end{tabular}
\end{table}

CL yields modest but significant improvements in ranking metrics: BEDROC +5.1\% ($p = 0.032$), EF@1\% +11.5\% ($p = 0.028$), Spearman $\rho$ +3.2\% ($p = 0.021$). These gains are consistent with the MSE improvements but smaller in magnitude, suggesting that CL primarily improves the accuracy of affinity predictions rather than their ranking.

\subsection{Cold-Target Split Results (GAT)}

\begin{table}[!h]
\caption{Ranking-level metrics on Davis cold-target split (GAT, 5 seeds).}
\label{tab:ranking_cold_target}
\centering
\begin{tabular}{lcccc}
\toprule
$\alpha$ & BEDROC($\alpha=20$) & EF@1\% & EF@5\% & Spearman $\rho$ \\
\midrule
0.000 (E0) & \textbf{0.312 $\pm$ 0.025} & \textbf{2.15 $\pm$ 0.32} & \textbf{1.42 $\pm$ 0.08} & \textbf{0.548 $\pm$ 0.019} \\
0.001 & 0.315 $\pm$ 0.024 & 2.18 $\pm$ 0.30 & 1.44 $\pm$ 0.07 & 0.552 $\pm$ 0.018 \\
0.010 & 0.305 $\pm$ 0.026 & 2.08 $\pm$ 0.33 & 1.39 $\pm$ 0.08 & 0.538 $\pm$ 0.020 \\
0.030 & 0.278 $\pm$ 0.029 & 1.85 $\pm$ 0.36 & 1.28 $\pm$ 0.09 & 0.502 $\pm$ 0.022 \\
0.050 & 0.258 $\pm$ 0.031 & 1.68 $\pm$ 0.38 & 1.20 $\pm$ 0.10 & 0.478 $\pm$ 0.024 \\
\bottomrule
\end{tabular}
\end{table}

Under cold-target splits, CL exhibits $\alpha$-dependent \textbf{ranking-level collapse}: at $\alpha = 0.03$, BEDROC degrades $-$10.9\% ($p = 0.003$), EF@1\% degrades $-$14.0\% ($p = 0.002$); at $\alpha = 0.05$, BEDROC degrades $-$17.3\% ($p < 0.001$), EF@1\% degrades $-$21.9\% ($p < 0.001$). Small $\alpha \leq 0.01$ causes negligible harm ($\leq 2.3\%$, $p > 0.15$). This confirms the $\alpha$-dependent vulnerability observed in MSE/CI also manifests in ranking metrics critical for virtual screening.

\subsection{Implications for Virtual Screening}

\begin{enumerate}
\item \textbf{Standard splits}: CL provides modest ranking improvements (+5--12\% in BEDROC/EF), suitable for retrospective virtual screening of compound libraries with known chemical space.

\item \textbf{Cold-target splits}: CL at $\alpha \geq 0.03$ causes significant ranking collapse ($-$10--22\% in BEDROC/EF), severely compromising early recognition of active compounds. For prospective screening against novel targets, $\alpha \leq 0.01$ or pure supervised training is strongly recommended.

\item \textbf{Practical recommendation}: If CL is used for virtual screening, validate on a project-specific cold-target or scaffold split before deployment. The standard-split ranking improvements do not transfer to genuinely novel targets or chemical space.
\end{enumerate}

\section{Summary}

This supplementary file provides:
\begin{enumerate}
\item \textbf{Validation-set $\alpha$ selection analysis}: Confirms Davis findings are robust ($-$0.2\% inflation), while KIBA estimates are conservative ($+6.9\%$ relative to test-set-selected).
\item \textbf{Ranking-level metrics}: Demonstrates CL improves standard-split ranking (+5--12\% BEDROC/EF) but causes cold-target ranking collapse ($-$10--22\% at $\alpha \geq 0.03$), with $\alpha \leq 0.01$ being safe.
\end{enumerate}

These analyses strengthen the main paper's conclusions and provide additional practical guidance for virtual screening applications.

\end{document}
